\documentclass[tikz,border=10pt]{standalone}
\usepackage{tikz}
\usepackage{amsmath}
\usepackage{amssymb}
\usepackage{ctex}
\usetikzlibrary{calc, positioning, arrows.meta, decorations.pathreplacing}

\begin{document}
\begin{tikzpicture}[>=Stealth, scale=1.0]

% ================================================================
% 步骤 1：u1 = v1
% ================================================================
\begin{scope}[xshift=0cm]
  \node[font=\small\bfseries, gray!70] at (1.5, 3.0) {步骤 1};

  % 坐标原点
  \coordinate (O1) at (0,0);
  \coordinate (V1) at (2.5, 1.2);

  % 背景淡网格
  \draw[gray!15] (-0.2,-0.2) grid (3.0,2.5);

  % v1 = u1（蓝色）
  \draw[->, thick, blue!80!black] (O1) -- (V1)
    node[above right, font=\small] {$\mathbf{u}_1 = \mathbf{v}_1$};

  \fill (O1) circle (1.5pt);

  % 标注框
  \node[font=\footnotesize, fill=blue!8, draw=blue!40, rounded corners=2pt,
        inner sep=4pt, align=center] at (1.5, -0.6)
    {$\mathbf{u}_1 = \mathbf{v}_1$};
\end{scope}

% ================================================================
% 步骤 2：u2 = v2 - proj_{u1} v2
% ================================================================
\begin{scope}[xshift=4.5cm]
  \node[font=\small\bfseries, gray!70] at (1.5, 3.0) {步骤 2};

  \coordinate (O2) at (0,0);
  \coordinate (U1) at (2.5, 0.5);   % u1 方向
  % v2 = (1.0, 2.2)
  \coordinate (V2) at (1.0, 2.2);
  % proj of v2 onto u1: (v2·u1)/(u1·u1)
  % v2·u1 = 2.5+0.55=3.05, u1·u1=6.25+0.25=6.5, scalar=3.05/6.5≈0.469
  % proj = 0.469*(2.5,0.5) = (1.173,0.235)
  \coordinate (P2) at (1.173, 0.235);
  % u2 = v2 - proj
  \coordinate (U2) at ($(V2)-(P2)$);  % (-0.173, 1.965)

  \draw[gray!15] (-0.2,-0.2) grid (3.0,2.5);

  % u1 方向（参考，淡灰）
  \draw[->, thick, gray!50] (O2) -- (U1)
    node[below right, font=\scriptsize, gray!70] {$\mathbf{u}_1$};

  % v2（橙色）
  \draw[->, thick, orange!80!black] (O2) -- (V2)
    node[left, font=\small] {$\mathbf{v}_2$};

  % 投影虚线（从 V2 垂直到 P2）
  \draw[densely dashed, gray!60] (V2) -- (P2);

  % 直角标记
  \draw[gray!60, thin]
    ($(P2)!0.18!(V2)$) --
    ($(P2)!0.18!(V2) + ($(P2)!0.18!(U1)$) - (P2)$) --
    ($(P2)!0.18!(U1)$);

  % proj_{u1}v2（红色虚线箭头）
  \draw[->, thick, red!60!black, dashed] (O2) -- (P2)
    node[below, font=\scriptsize, yshift=-3pt]
      {$\mathrm{proj}_{\mathbf{u}_1}\mathbf{v}_2$};

  % u2（绿色）
  \draw[->, thick, green!60!black] (O2) -- (U2)
    node[left, font=\small] {$\mathbf{u}_2$};

  \fill (O2) circle (1.5pt);

  \node[font=\footnotesize, fill=green!8, draw=green!50, rounded corners=2pt,
        inner sep=4pt, align=center] at (1.5, -0.6)
    {$\mathbf{u}_2 = \mathbf{v}_2 - \mathrm{proj}_{\mathbf{u}_1}\mathbf{v}_2$};
\end{scope}

% ================================================================
% 步骤 3：u3 = v3 - proj_{u1}v3 - proj_{u2}v3
% ================================================================
\begin{scope}[xshift=9.0cm]
  \node[font=\small\bfseries, gray!70] at (1.5, 3.0) {步骤 3};

  \coordinate (O3) at (0,0);
  % u1, u2 方向（正交）
  \coordinate (U1b) at (2.2, 0.0);
  \coordinate (U2b) at (0.0, 2.0);
  % v3 = (1.6, 1.5)
  \coordinate (V3) at (1.6, 1.5);
  % proj onto u1: (1.6/2.2)*(2.2,0) = (1.6, 0)
  \coordinate (P3u1) at (1.6, 0.0);
  % proj onto u2: (1.5/2.0)*(0,2.0) = (0, 1.5)
  \coordinate (P3u2) at (0.0, 1.5);
  % u3 = v3 - P3u1 - P3u2 = (0, 0) for this demo
  % 用稍微不整齐的数值让图示更真实
  % v3 = (1.8, 1.6), u1=(2.5,0.5)norm, u2=(0,1)norm
  % 重新设置更直观的值
  % 让 v3=(1.5,1.5), u1=(2.2,0), u2=(0,2), proj1=(1.5,0), proj2=(0,1.5), u3=(0,0)
  % 换成 v3=(1.5,1.7) => u3=(0,0.2) 更好看
  \coordinate (V3) at (1.5, 1.7);
  \coordinate (P3u1) at (1.5, 0.0);
  \coordinate (P3u2) at (0.0, 1.5);
  \coordinate (U3) at (0.0, 0.2);

  \draw[gray!15] (-0.2,-0.2) grid (3.0,2.5);

  % u1（淡）
  \draw[->, thick, gray!40] (O3) -- (U1b)
    node[below right, font=\scriptsize, gray!60] {$\mathbf{u}_1$};
  % u2（淡）
  \draw[->, thick, gray!40] (O3) -- (U2b)
    node[left, font=\scriptsize, gray!60] {$\mathbf{u}_2$};

  % v3（橙色）
  \draw[->, thick, orange!80!black] (O3) -- (V3)
    node[right, font=\small] {$\mathbf{v}_3$};

  % proj u1 虚线
  \draw[->, thick, red!50!black, dashed] (O3) -- (P3u1)
    node[below, font=\scriptsize] {$\mathrm{proj}_{\mathbf{u}_1}\mathbf{v}_3$};

  % proj u2 虚线
  \draw[->, thick, purple!60!black, dashed] (O3) -- (P3u2)
    node[left, font=\scriptsize] {$\mathrm{proj}_{\mathbf{u}_2}\mathbf{v}_3$};

  % 垂直辅助线从 v3 到投影点
  \draw[densely dashed, gray!50] (V3) -- (P3u1);
  \draw[densely dashed, gray!50] (V3) -- (P3u2);

  % u3（绿色）
  \draw[->, very thick, green!60!black] (O3) -- (U3)
    node[right, font=\small, xshift=4pt] {$\mathbf{u}_3$};

  \fill (O3) circle (1.5pt);

  \node[font=\footnotesize, fill=green!8, draw=green!50, rounded corners=2pt,
        inner sep=4pt, align=center] at (1.5, -0.6)
    {$\mathbf{u}_3 = \mathbf{v}_3 - \mathrm{proj}_{\mathbf{u}_1}\mathbf{v}_3 - \mathrm{proj}_{\mathbf{u}_2}\mathbf{v}_3$};
\end{scope}

% ================================================================
% 顶部标题
% ================================================================
\node[font=\large\bfseries] at (6.5, 3.9) {Gram-Schmidt 正交化过程};

% 底部公式框（整体）
\node[font=\small, fill=white, draw=gray!50, thin, rounded corners=3pt,
      inner sep=7pt, align=center] at (6.5, -1.6) {
  $\mathbf{u}_1 = \mathbf{v}_1$，\quad
  $\mathbf{u}_k = \mathbf{v}_k - \displaystyle\sum_{j=1}^{k-1}
    \frac{\langle\mathbf{v}_k,\mathbf{u}_j\rangle}{\langle\mathbf{u}_j,\mathbf{u}_j\rangle}\mathbf{u}_j$
  \quad ($k=2,3,\ldots$)
};

\end{tikzpicture}
\end{document}
